%NO MODIFICAR ESTA SECCION!
\documentclass[12pt]{article} % Define la clase del documento, en este caso, un artículo
\usepackage{mathptmx}
\usepackage[letterpaper,margin=3cm]{geometry} % Configura el tamaño del papel y los márgenes del documento
\usepackage{graphicx} % Permite la inserción de imágenes
\usepackage[spanish]{babel}% Activar esta configuración para informes en español, ajusta el idioma del documento
\usepackage[usenames]{color} % Permite el uso de colores definidos por nombre en el documento
\usepackage{hyperref} % Habilita enlaces y referencias dentro del documento
\usepackage{booktabs} % Proporciona comandos para crear tablas de alta calidad
\usepackage{natbib}
\usepackage{tikz} % Permite la creación de gráficos y diagramas vectoriales directamente en LaTeX
\usepackage{float} % Para controlar la posición de los elementos flotantes, como imágenes, con la opción [H]
\usepackage{diagbox} % Permite crear celdas con líneas diagonales en tablas
\usepackage{listings} % Permite la inclusión y formateo de código fuente en el documento
\usepackage{xcolor} % Paquete para definir y usar colores en el documento
\usepackage{parskip} % Añade espacio entre párrafos en lugar de sangrías
\usepackage{fancyhdr} % Permite personalizar encabezados y pies de página
\usepackage{amsmath} % Proporciona una amplia variedad de entornos y comandos matemáticos
\usepackage{microtype}  % Mejora la justificación del texto
\usepackage{ragged2e}


\hypersetup{
    colorlinks   = true,    % Colorea los enlaces
    linkcolor    = blue,    % Color de los enlaces internos (por ejemplo, en la tabla de contenidos)
    citecolor    = black,    % Color de los enlaces a citas bibliográficas
    filecolor    = blue,    % Color de los enlaces a archivos locales
    urlcolor     = blue,    % Color de los enlaces externos (URLs)
}


\geometry{
    paperwidth=21.6cm,  % Ancho del papel
    paperheight=27.9cm,  % Largo del papel
    left=3cm,  % Margen izquierdo
    right=2cm,  % Margen derecho
}

\pagestyle{fancy} % Usa el estilo fancyhdr
\fancyhf{} % Borra todos los encabezados y pies de página
\renewcommand{\headrulewidth}{0pt}
\renewcommand{\footrulewidth}{0pt} % Desactiva la línea horizontal predeterminada en el pie
%\setlength{\headheight}{1.5cm} % Ajusta la altura del encabezado para hacer espacio para la línea
%\fancyhead[L]{\raisebox{0.20cm}{\textbf{Sistemas de Transporte}}} % Añade el texto en la parte izquierda del encabezado, subiéndolo ligeramente
%\fancyhead[R]{\raisebox{0.1cm}{\includegraphics[width=0.25\linewidth]{LOGO_UNIVERSIDAD.jpg}}} % Añade la imagen en la parte derecha del encabezado y súbela un poco
%\fancyhead[C]{\rule{\textwidth}{0.6pt}} % Añade una línea horizontal superior centrada
%\fancyfoot[C]{\rule{\textwidth}{0.6pt}} % Añade una línea horizontal en el pie de página centrada
%\fancyfoot[R]{\raisebox{-1.5\baselineskip}{\thepage}} % Coloca el número de página a la derecha, con suficiente espacio debajo de la línea
%\geometry{top=2.5cm, bottom=2.5cm} % Ajusta los márgenes superior e inferior

\fancyhead[R]{\thepage} % Número de página a la derecha en el encabezado


% Definición de colores al estilo Visual Studio Code
\definecolor{codegreen}{rgb}{0.25,0.49,0.48} % Comentarios
\definecolor{codegray}{rgb}{0.5,0.5,0.5} % Números y anotaciones
\definecolor{codepurple}{rgb}{0.58,0,0.82} % Palabras clave
\definecolor{backcolour}{rgb}{0.95,0.95,0.92} % Color de fondo

% Configuración del estilo de las celdas de código
\lstset{
    backgroundcolor=\color{backcolour},   % color de fondo; necesita que el paquete color o xcolor esté cargado
    commentstyle=\color{codegreen},       % estilo de comentarios
    keywordstyle=\color{codepurple},      % estilo de palabras clave
    numberstyle=\tiny\color{codegray},    % estilo de los números de línea
    stringstyle=\color{red},              % estilo de las cadenas de texto
    basicstyle=\ttfamily\small,           % estilo del texto básico
    breakatwhitespace=false,              % ajustes de líneas sólo en espacios en blanco
    breaklines=true,                      % ajustar las líneas si son muy largas
    captionpos=b,                         % posición de la leyenda (abajo)
    keepspaces=true,                      % preserva los espacios en el texto; útil si se usa monoespaciado
    numbers=left,                         % dónde poner los números de línea
    numbersep=5pt,                        % qué tan lejos están los números de línea del código
    showspaces=false,                     % mostrar espacios con subrayados particulares; reemplaza 'showstringspaces'
    showstringspaces=false,               % subrayar los espacios dentro de las cadenas solo
    showtabs=false,                       % mostrar tabulaciones en el código con subrayados particulares
    tabsize=2,                            % tamaños de tabulación a 2 espacios
    language=TeX,                         % lenguaje del código
    morecomment=[l]\#,                    % reconocer # como inicio de comentario en Python
    frame=single,                         % agregar un marco simple alrededor del código
    rulecolor=\color{black}               % color del marco
}

\begin{document}
%----------------------------------------------------------------------------------------
%   PORTADA
%Modificar desde aqui en adelante
%----------------------------------------------------------------------------------------
\begin{titlepage}%Inicio de la carátula, solo modificar los datos necesarios
\newcommand{\HRule}{\rule{\linewidth}{0.5mm}} 
\center 
%----------------------------------------------------------------------------------------
%	ENCABEZADO
%----------------------------------------------------------------------------------------
%----------------------------------------------------------------------------------------
%	SECCION DEL TITULO
%----------------------------------------------------------------------------------------
\begin{center}
    \textbf{\LARGE Generación y Atracción de Viajes} \\[0.5cm]
    \textbf{Felipe Vicencio y Lukas Wolff} \\
    Facultad de Ingenieria y Ciencias Aplicadas, Universidad de los Andes, Santiago de Chile.\\
    email: \href{mailto:lwolff@miuandes.cl}{lwolff@miuandes.cl}, \href{mailto:favicencio@miuandes.cl}{favicencio@miuandes.cl}
    \\
    github: \href{https://github.com/LukasWolff2002/TAREA_2_AUTITOS.git}{Link al repositorio}
\end{center}

\vspace{1cm}

\begin{center}
    \textbf{\large ABSTRACT}    
\end{center}

\begin{justify}
    Nowadays, the importance of having an efficient \textit{transport 
    system} and a precise analysis of travels per person to 
    determine the flux and social behavior is key to the 
    well-functioning of the country.  In this paper, the commune 
    of \textit{Las Condes} was analyzed with a \textit{linear regression} equation to 
    determine the number of \textit{travels} per home in each of the years 2012, 
    2017 and 2023. The variables were separated according to age range 
    and per capita income. It is worth highlighting that the analysis 
    was made to compare a sample data (730 homes) from the \textit{ESI survey} (\textbf{\cite{esi}}), with the population data (321,110 homes)
    from a data base generated by the \textit{2017 census} (\textbf{\cite{censo2017}}). 
    The results show that the average amount of people living in a house has decreased 
    in the last 11 years from 3.16 to 2.6 people per household. Also, the number of trips 
    made in a household did vary between 2012 and 2023 in Las Condes, decreasing form 
    1.84 to 1.66 trips. Although, if we took into account the variability of the inflation, IPC, gas, dollar and public transport price, 
    and the pandemic, the number of trips throughout the years increased in almost 25\%, due to the increase in the number of houses. 
    \\ \\ \\ \\ \\
    \textbf{Key Words}: Linear Regression; Travels; Transport System; Las Condes; 2017 Census; ESI Surevey.
\end{justify}


\vspace{1cm}

\end{titlepage}
%----------------------------------------------------------------------------------------
%	SECCION DE LA FECHA
%----------------------------------------------------------------------------------------
%{\large \textbf{\today}}\\[2cm] % El comando \today coloca la fecha del dia, y esto se actualiza con cada compilacion, en caso de querer tener una fecha estatica, reemplazar el \today por la fecha deseada

%----------------------------------------------------------------------------------------
%  INDICE
%----------------------------------------------------------------------------------------
%\newpage
%\tableofcontents
%\thispagestyle{plain} % Deshabilita el encabezado en la página del índice
%\thispagestyle{empty} % Deshabilita el número de página en la página del índice
%\newpage

%Se puede agregar un indice de figuras si es nesesario
%\newpage
%\listoffigures 
%\thispagestyle{plain} % Deshabilita el encabezado en la página del índice %
%\thispagestyle{empty}
%\newpage
%----------------------------------------------------------------------------------------
%   ACÁ EMPIEZA EL INFORME
\setcounter{page}{1} % Reinicia el contador de páginas
%----------------------------------------------------------------------------------------
%Este es el formato a seguir para los titulos de las secciones

\section{Introducción}

El estudio de los viajes es fundamental para el desarrollo de zonas altamente pobladas, ya que permite ahorrar tiempo de viaje, reducir el consumo de combustible y mejorar el bienestar de las personas. Para generar estos análisis, es necesario recolectar datos reales sobre el comportamiento de las personas, como su tipo de transporte, origen y destino, horas de viaje, entre otros.
\\ \\
El problema de estos estudios radica en la cantidad de datos que se necesitan para realizar un análisis completo. Por lo tanto, se recurre a estimaciones mediante una regresión lineal para representar a la población de la manera más eficaz posible.
\\ \\
En este informe, se buscará analizar el número de viajes de la comuna de Las Condes, basándose en los datos obtenidos de la encuesta ESI para los años 2012, 2017 y 2023. Se utilizará un modelo de regresión lineal para estimar el número de viajes de la comuna en base a distintos factores que pueden influir en la cantidad de viajes.
\\ \\
Finalmente, se analizará y discutirá la eficiencia del modelo, la precisión de los datos obtenidos y los factores que pueden influir en la cantidad de viajes en la comuna de Las Condes.

\section{Contenido}

\subsection{Matriz Origen-Destino}

Para representar todos los viajes de la red, se puede utilizar la matriz origen-destino, la cual se representa de la siguiente manera:

\begin{equation}
    \begin{bmatrix}
        O_{1}D_{1} & O_{1}D_{2} & \cdots & O_{1}D_{n} \\
        O_{2}D_{1} & O_{2}D_{2} & \cdots & O_{2}D_{n} \\
        \vdots & \vdots & \ddots & \vdots \\
        O_{n}D_{1} & O_{n}D_{2} & \cdots & O_{n}D_{n}
    \end{bmatrix}
\end{equation}

Generalmente no se disponen de suficientes datos para completar satisfactoriamente la matriz origen-destino, por lo que se deben utilizar modelos de estimación de viajes para completarla.
\\ \\
Es importante analizar los distintos factores que pueden influir en los viajes de las personas.

\subsection{Regresión Lineal}

Basándonos en los datos obtenidos de distintas encuestas, se busca establecer una relación que explique el comportamiento de la población:

\begin{equation}
    Y = \alpha + \beta_{1}X_{1} + \beta_{2}X_{2} + \cdots + \beta_{n}X_{n} + \epsilon
\end{equation}

Que se puede extrapolar a:

\begin{equation}
    Y = \alpha + \vec{\beta}\vec{X} + \epsilon
\end{equation}

Donde $Y$ corresponde al número de viajes y $\vec{X}$ a los distintos factores que pueden afectar la cantidad de viajes. En este informe, se utilizarán las siguientes variables $X$:

\begin{center}
    $X_i = \log_{10}(IPCH)$, donde IPCH corresponde al ingreso per cápita por hogar.\\
    $X_p$ = número de personas con edad $\epsilon$ [0, 5].\\
    $X_e$ = número de personas con edad $\epsilon$ [6, 22].\\
    $X_t$ = número de personas con edad $\epsilon$ [23, 62].\\
    $X_j$ = número de personas con edad $\epsilon$ [63, 79].
\end{center}

De esta forma, el modelo de regresión lineal queda de la siguiente manera:

\begin{equation}
    Y_i = -2.1723 + 0.3792X_i + 0.6221X_p + 1.0065X_e + 0.4302X_t + 0.1614X_j
    \label{eq:regresion}
\end{equation}

Es importante mencionar que los coeficientes $\vec{\beta}$ y $\alpha$ fueron proporcionados por el enunciado, donde el modelo fue calibrado según los datos del 2012. Además, se asume que este modelo cumple los tres supuestos correspondientes: independencia, linealidad y homocedasticidad.

\section{Obtención de Resultados a partir de Bases de Datos}

Para el análisis y discusión de la regresión lineal, se utilizará la base de datos ESI \textbf{\cite{esi}}, la cual contiene toda la información necesaria para aplicar en el modelo.
\\ \\
Inicialmente, se realizó un filtrado de datos correspondientes a la comuna de Las Condes, además de calcular los distintos factores y promediarlos para los años 2012, 2017 y 2023.

\begin{table}[H]
    \centering
    \caption{Coeficientes promediados}
    \vspace{0.2cm}
    \begin{tabular}{|c|c|c|c|c|c|c|c|}
        \hline
        Año & Número Hogares & $X_i$ & $X_p$ & $X_e$ & $X_t$ & $X_j$ & Tamaño hogar \\
        \hline
        2012 & 271 & 5.730 & 0.258 & 0.771 & 1.978 & 0.303 & 3.16\\
        2017 & 272 & 5.773 & 0.243 & 0.691 & 1.978 & 0.320 & 3.09\\
        2023 & 187 & 5.934 & 0.160 & 0.668 & 1.775 & 0.299 & 2.60\\
        \hline
    \end{tabular}
    \vspace{0.2cm}
    \\Fuente: Elaboración propia a partir de los datos de encuesta ESI \textbf{\cite{esi}}
\end{table}

Posteriormente, se consultó la encuesta de Estimación y Proyección, de esta forma se tendrán los datos reales sobre la población, para poder comparar con la muestra de la encuesta ESI.

\begin{table}[H]
    \centering
    \caption{Valores Población}
    \vspace{0.2cm}
    \begin{tabular}{|c|c|c|c|c|c|}
        \hline
        \diagbox{Año}{Coeficiente} & $X_p$ & $X_e$ & $X_t$ & $X_j$ & Total hogares \\
        \hline
        2012 & 0.219 & 0.705 & 1.751 & 0.483 & 90304 \\
        2017 & 0.204 & 0.611 & 1.736 & 0.537 & 99582 \\
        2023 & 0.156 & 0.457 & 1.475 & 0.510 & 131224 \\
        \hline
    \end{tabular}
    \vspace{0.2cm}
    \\Fuente: Elaboración propia a partir de los datos de Estimación y Proyección 2002 - 2035
\end{table}

Conociendo el total de hogares de la población, es posible calcular el número de viajes totales obtenidos a partir de la regresión lineal:

\begin{table}[H]
    \centering
    \caption{Estimación de viajes totales}
    \vspace{0.2cm}
    \begin{tabular}{|c|c|c|}
        \hline
        Año & $Y_i$ & Viajes Totales \\
        \hline
        2012 & 1.84 & 192554\\
        2017 & 1.77 & 208872\\
        2023 & 1.66 & 195891\\
        \hline
    \end{tabular}
    \vspace{0.2cm}
    \\ Fuente: Elaboración propia a partir de los resultados de la ecuación \ref{eq:regrecion}
\end{table}

*Los viajes totales se truncaron hacia abajo, ya que el número de viajes no puede ser decimal. De esta forma, no es posible generar un viaje extra.

\section{Discusiones}

\subsection{Inflación}

Un punto clave a considerar sobre el comportamiento de las personas es el valor y peso que tiene la moneda en el tiempo. Por lo tanto, es necesario considerar la inflación y el aumento del precio en los distintos insumos del transporte.

\begin{table}[H]
    \centering
    \caption{Inflación en Chile}
    \vspace{0.2cm}
    \begin{tabular}{|c|c|}
        \hline
        Periodo & Inflación (\%) \\
        \hline
        2012 - 2017 & 18,2 \\
        2012 - 2023 & 61,5 \\
        \hline
    \end{tabular}
    \vspace{0.2cm}
    \\Fuente: \textbf{\cite{ipc}}
\end{table}

La regresión lineal aplicada fue calibrada según los datos de 2012, por lo tanto, se puede ajustar $X_i$ según la inflación para obtener un parámetro de viajes más acertado.

\begin{table}[H]
    \centering
    \caption{Ajuste de modelo por inflación}
    \vspace{0.2cm}
    \begin{tabular}{|c|c|c|c|}
        \hline
        Año & $X_i$ ajustado & $Y_i$ & Viajes Totales \\
        \hline
        2012 & 5.730 & 1.84 & 166159 \\
        2017 & 5.700 & 1.74 & 173272 \\
        2023 & 5.725 & 1.58 & 207333 \\
        \hline
    \end{tabular}
    \vspace{0.2cm}
    \\Fuente: Elaboración propia
\end{table}

*Los viajes totales se truncaron hacia abajo, ya que el número de viajes no puede ser decimal, evitando así la generación de viajes extra.

\subsection{Precisión de Resultados}

Analizando los datos presentados junto con las bases de datos, se estima que la precisión de los resultados obtenidos es fiable, ya que las fuentes de donde se recopiló la información provienen directamente de páginas web gubernamentales, municipales y de agentes encargados de análisis estadísticos sociales, como el INE.

Si bien se evidencia una variabilidad en los parámetros $X_p$, $X_e$, $X_t$ y $X_j$, esto se debe a la comparación de datos muestrales con datos poblacionales. Estos datos fueron comparados estadísticamente, como se muestra en la Tabla \textbf{\ref{Precision_Resultados}} (tabla no incluida).

Se puede apreciar que los resultados poseen una desviación estándar insignificante, lo que indica que son representativos para cada coeficiente. Además, se realizó un análisis con el estadístico t para comparar la media muestral con la media poblacional. Debido a que \(|p_{value}| < |t_{value}|\), existe una correlación significativa entre las medias, lo que permite afirmar que los resultados son precisos y descriptivos.

\subsection{Análisis}

Considerando todo lo expuesto anteriormente, se puede observar que el modelo de regresión lineal, sin ser ajustado, es fiable con cierto margen de error. Se concluye que el número de viajes disminuye por cada casa, lo cual tiene sentido al analizar la inflación y el valor de los distintos insumos de transporte (ver Tabla \textbf{\ref{Insumos}}). Además, se puede observar que el tamaño de las familias disminuye, lo cual influye en el número de viajes. Finalmente, es necesario considerar el efecto que tuvo la pandemia sobre las personas, donde muchos trabajos comenzaron a ser virtuales o semi presenciales, lo cual afecta el transporte en familias, sobre todo las de menor tamaño.
\\ \\
Se observa que al calibrar el modelo según la inflación, los viajes por hogar bajan aún más, lo cual se relaciona con la pérdida de valor que tiene el dinero y no ha sido proporcional a los aumentos de sueldos (ver Tabla \ref{VariacionIPCH}).
\\ \\
Finalmente, haciendo un análisis estadístico de los datos, se concluye que los datos se pueden considerar representativos y precisos, dada la procedencia de estos y su correlación muestral con poblacional según el estadístico T Student.

\section{Conclusión}

En base a los resultados obtenidos, se puede concluir que el modelo de regresión lineal es fiable para estimar el número de viajes de la comuna de Las Condes; sin embargo, es necesario ajustar el modelo según la inflación para obtener resultados más precisos.
\\ \\
El número de viajes por hogar disminuye año a año, lo cual se explica por la disminución del tamaño familiar, la falta de relación entre la pérdida de valor de la moneda y el aumento del sueldo, además del efecto que tuvo la pandemia en el comportamiento de las personas.
\\ \\
Considerando todo lo anterior, se puede concluir que todos los objetivos propuestos fueron alcanzados, por lo que la práctica fue completada de manera satisfactoria.

\section{Anexos}

\subsection{Tablas Inflación}

\begin{table}[H]
    \centering
    \caption{Aumento de indicadores económicos}
    \vspace{0.2cm}
    \begin{tabular}{|c|c|c|c|c|}
        \hline
        Año & UF (Pesos) & $\Delta$ \% UF & Valor Dólar & $\Delta$ \% Dólar \\
        \hline
        2012 & 22,296.19 & - & 501.34 & - \\
        2017 & 26,348.83 & 18.2\% & 661.19 & 31.88\% \\
        2023 & 35,122.26 & 57.6\% & 826.34 & 64.82\% \\
        \hline
    \end{tabular}
    \vspace{0.2cm}
    \\Fuente: \textbf{\cite{sii}}
\end{table}

\begin{table}[H]
    \centering
    \caption{Aumento de precio de insumos}
    \vspace{0.2cm}
    \begin{tabular}{|c|c|c|c|c|}
        \hline
        Año & Bencina 95 (USD) & $\Delta$ \% Bencina & Metro (CLP) & $\Delta$ \% Metro \\
        \hline
        2012 & 1.56 & - & 610 & - \\
        2017 & 1.15 & -26.28\% & 660 & 8.19\% \\
        2023 & 1.64 & 5.13\% & 730 & 19.67\% \\
        \hline
    \end{tabular}
    \label{Insumos}
    \vspace{0.2cm}
    \\Fuente: \textbf{\cite{tradingeconomics}, \cite{bcentral}}
\end{table}

\begin{table}[H]
    \centering
    \caption{Variación IPCH ajustado por inflación}
    \vspace{0.2cm}
    \begin{tabular}{|c|c|c|}
        \hline
        Año & $X_i$ & $\Delta$ \% \\
        \hline
        2012 & 5.730 & - \\
        2017 & 5.700 & -0.52\% \\
        2023 & 5.725 & -0.09\% \\
        \hline
    \end{tabular}
    \vspace{0.2cm}
    \label{VariacionIPCH}
    \\Fuente: Elaboración propia
\end{table}

\subsection{Precisión de Datos} \label{sec:Precision}

\begin{table}[H]
    \caption{Análisis estadístico para la precisión de resultados.}
    \vspace{0.2cm}
    \centering
    \begin{tabular}{|c|c|c|c|}
    \hline
    \textbf{Coeficiente} & \textbf{Estadístico} & \textbf{Población} & \textbf{Muestra} \\ \hline
    $X_p$ & \multicolumn{1}{|c|}{\(\mu\)} & 0.193 & 0.220 \\ \cline{2-4}
    & \multicolumn{1}{|c|}{\(\sigma\)} & 0.027 & 0.053 \\ \cline{2-4}
    & \multicolumn{1}{|c|}{Valor p} & \multicolumn{2}{|c|}{1} \\ \cline{2-4}
    & \multicolumn{1}{|c|}{Valor t} & \multicolumn{2}{|c|}{-26.09} \\ \hline
    
    $X_e$ & \multicolumn{1}{|c|}{\(\mu\)} & 0.591 & 0.710 \\ \cline{2-4}
    & \multicolumn{1}{|c|}{\(\sigma\)} & 0.102 & 0.054 \\ \cline{2-4}
    & \multicolumn{1}{|c|}{Valor p} & \multicolumn{2}{|c|}{1} \\ \cline{2-4}
    & \multicolumn{1}{|c|}{Valor t} & \multicolumn{2}{|c|}{-31.52} \\ \hline
    
    $X_t$ & \multicolumn{1}{|c|}{\(\mu\)} & 1.654 & 1.910 \\ \cline{2-4}
    & \multicolumn{1}{|c|}{\(\sigma\)} & 0.126 & 0.117 \\ \cline{2-4}
    & \multicolumn{1}{|c|}{Valor p} & \multicolumn{2}{|c|}{1} \\ \cline{2-4}
    & \multicolumn{1}{|c|}{Valor t} & \multicolumn{2}{|c|}{-55.84} \\ \hline
    
    $X_j$ & \multicolumn{1}{|c|}{\(\mu\)} & 0.510 & 0.307 \\ \cline{2-4}
    & \multicolumn{1}{|c|}{\(\sigma\)} & 0.021 & 0.011 \\ \cline{2-4}
    & \multicolumn{1}{|c|}{Valor p} & \multicolumn{2}{|c|}{0} \\ \cline{2-4}
    & \multicolumn{1}{|c|}{Valor t} & \multicolumn{2}{|c|}{261.09} \\ \hline
    \end{tabular}
    \label{Precision_Resultados}
    \vspace{0.2cm}
    \\Fuente: Elaboración propia.
\end{table}

\subsection{Tablas Datos Encuesta EP} \label{sec:Datos_EP}

\begin{table}[H]
    \centering
    \caption{Rango etario en Las Condes año 2012}
    \vspace{0.2cm}
    \begin{tabular}{|c|c|c|c|c|c|}
        \hline
        \textbf{Rango etario} & \textbf{Hombres} & \textbf{Mujeres} & \textbf{Total} & \textbf{\% Hombres} & \textbf{\% Mujeres} \\ \hline
        0-5 & 10484 & 9377 & 19861 & 52.79\% & 47.21\% \\ \hline
        6-22 & 33329 & 30366 & 63695 & 52.32\% & 47.68\% \\ \hline
        23-62 & 75938 & 82177 & 158115 & 48.02\% & 51.98\% \\ \hline
        63-80 & 17595 & 26095 & 43690 & 40.27\% & 59.73\% \\ \hline
    \end{tabular}
    \label{Cuadro 4}
    \vspace{0.2cm}
    \\Fuente: Elaboración propia.
\end{table}

\begin{table}[H]
    \centering
    \caption{Rango etario en Las Condes año 2017}
    \vspace{0.2cm}
    \begin{tabular}{|c|c|c|c|c|c|}
        \hline
        \textbf{Rango etario} & \textbf{Hombres} & \textbf{Mujeres} & \textbf{Total} & \textbf{\% Hombres} & \textbf{\% Mujeres} \\ \hline
        0-5 & 10672 & 9692 & 20364 & 52.41\% & 47.59\% \\ \hline
        6-22 & 31833 & 29028 & 60861 & 52.31\% & 47.69\% \\ \hline
        23-62 & 84879 & 88035 & 172914 & 49.08\% & 50.92\% \\ \hline
        63-80 & 21753 & 31816 & 53569 & 40.61\% & 59.39\% \\ \hline
    \end{tabular}
    \label{Cuadro 5}
    \vspace{0.2cm}
    \\Fuente: Elaboración propia.
\end{table}

\begin{table}[H]
    \centering
    \caption{Rango etario en Las Condes año 2023}
    \vspace{0.2cm}
    \begin{tabular}{|c|c|c|c|c|c|}
        \hline
        \textbf{Rango etario} & \textbf{Hombres} & \textbf{Mujeres} & \textbf{Total} & \textbf{\% Hombres} & \textbf{\% Mujeres} \\ \hline
        0-5 & 10769 & 9758 & 20527 & 52.46\% & 47.54\% \\ \hline
        6-22 & 31416 & 28678 & 60094 & 52.27\% & 47.73\% \\ \hline
        23-62 & 96289 & 97341 & 193630 & 49.72\% & 50.28\% \\ \hline
        63-80 & 27825 & 39107 & 66932 & 41.57\% & 58.43\% \\ \hline
    \end{tabular}
    \label{Cuadro 6}
    \vspace{0.2cm}
    \\Fuente: Elaboración propia.
\end{table}

\bibliographystyle{plainnat}  % Estilo de la bibliografía (puedes cambiarlo a otro si prefieres)
\bibliography{referencias.bib}  % Nombre del archivo .bib

\end{document}
